\documentclass[UTF8,12pt,a4paper]{ctexart}

\usepackage{amsmath}
\usepackage{amssymb}
\usepackage{graphicx}
\usepackage{enumerate}
\usepackage[margin=1in]{geometry}
\geometry{a4paper}
\usepackage{fancyhdr}
\pagestyle{fancy}
\fancyhf{}
\usepackage{xcolor}
\usepackage{hyperref}
\hypersetup{
    pdfborder={0 0 0},
    colorlinks=true,
    linkcolor=black,
    citecolor=black,
    urlcolor=blue
}
\usepackage{booktabs}
\usepackage{array}
\usepackage{multirow}
\usepackage{longtable}
\newcolumntype{L}[1]{>{\raggedright\arraybackslash}p{#1}}
\newcolumntype{M}[1]{>{\centering\arraybackslash}p{#1}}

\title{集成电路工艺技术基础 Chapter 7\\光刻工艺详细复习手稿}
\author{冯峻}
\date{\today}

\fancyhead[L]{Chapter 7 复习手稿}
\fancyhead[C]{光刻工艺}
\fancyfoot[C]{\thepage}

\begin{document}

\maketitle
\tableofcontents
\newpage

\section{光刻技术基本概念}

\subsection{光刻的起源与发展}

\subsubsection{传统光刻技术的发明}

\textbf{发明者：}Alois Senefelder

\textbf{发明时间：}1796年

\textbf{原始方法：}化学印刷（Chemical Printing）

\textbf{工艺原理：}
\begin{enumerate}
    \item 用油基物质在石灰石上标记图案
    \item 酸将图案蚀刻到表面
    \item 阿拉伯树胶溶液仅粘附于非油性区域
    \item 印刷时，水附着在树胶区域，油性墨（肥皂、蜡、油、煤黑）转移到纸上
\end{enumerate}

\textbf{词源：}
\begin{itemize}
    \item lithos = 石头（希腊语）
    \item graphy = 书写
\end{itemize}

\subsubsection{现代半导体光刻}

\textbf{在IC制造中的地位：}
\begin{itemize}
    \item 平面工艺中的关键工序
    \item 成本占比超过60\%
    \item 决定器件的最小尺寸和性能
\end{itemize}

\textbf{当前生产技术（2020年代）：}
\begin{itemize}
    \item \textbf{先进节点：}5 nm EUV光刻（TSMC）
    \item \textbf{曝光波长：}193 nm（ArF浸没式）或13.5 nm（EUV）
    \item \textbf{分辨率：}约$\lambda$/4，可达$\sim$5 nm（使用分辨率增强技术）
    \item \textbf{产能：}100片晶圆/小时（300 mm晶圆）
    \item \textbf{光刻胶：}化学放大型光刻胶（Chemically Amplified Resist）
\end{itemize}

\subsubsection{极限研究光刻}

\textbf{STM原子操纵技术：}
\begin{itemize}
    \item \textbf{方法：}扫描隧道显微镜（STM）逐个拾取和放置原子
    \item \textbf{条件：}
    \begin{itemize}
        \item 超高真空（UHV）
        \item 低温环境
        \item 原子级清洁表面
        \item 原子级尖锐的STM针尖
    \end{itemize}
    \item \textbf{分辨率：}$\sim$0.3 nm（单原子）
    \item \textbf{产能：}1原子/分钟$\approx$0.017像素/秒
    \item \textbf{局限性：}非通用工艺，仅用于研究
\end{itemize}

\subsection{光刻技术的关键指标}

\subsubsection{分辨率（Resolution）}

\textbf{定义：}能够清晰成像的最小特征尺寸

\textbf{衡量标准：}
\begin{itemize}
    \item 最小线宽
    \item 尺寸控制精度（Dimensional control）
    \begin{itemize}
        \item 跨样品（across sample）
        \item 样品间（sample to sample）
    \end{itemize}
\end{itemize}

\subsubsection{对比度（Contrast）}

\textbf{定义：}区分曝光和未曝光区域的能力

\textbf{重要性：}直接影响图形边缘的清晰度

\subsubsection{产能（Throughput）}

\textbf{定义：}单位时间内完成曝光的晶圆数量

\textbf{影响因素：}
\begin{itemize}
    \item 像素曝光速率
    \item 闪光速率（Flash rate）
    \item 每次闪光的像素数
    \item 光刻胶灵敏度
\end{itemize}

\subsubsection{套刻精度（Placement Accuracy/Overlay）}

\textbf{定义：}不同层之间图形对准的精度

\textbf{类型：}
\begin{itemize}
    \item 层间套刻（Layer to layer）
    \item 特征间套刻（Feature to feature）
\end{itemize}

\subsection{通用光刻工艺流程}

\textbf{光刻作为信息流动过程：}

设计图形$\rightarrow$掩模制作$\rightarrow$曝光转移$\rightarrow$显影$\rightarrow$图形转移

\textbf{完整光刻工艺步骤：}
\begin{enumerate}
    \item \textbf{涂胶（Spincoat）：}旋涂光刻胶
    \item \textbf{前烘（Post-apply bake）：}溶剂挥发
    \item \textbf{曝光（Exposure）：}紫外光通过掩模曝光
    \item \textbf{曝光后烘焙（Post-exposure bake, PEB）：}促进化学反应
    \item \textbf{显影（Development）：}溶解可溶部分
    \item \textbf{干燥（Drying）：}去除显影液
\end{enumerate}

\section{光刻技术与曝光系统}

\subsection{光源（Sources）}

\subsubsection{汞灯（Hg Arc Lamp）}

\textbf{主要谱线：}
\begin{itemize}
    \item \textbf{G线：}436 nm
    \item \textbf{H线：}405 nm
    \item \textbf{I线：}365 nm（最常用）
\end{itemize}

\textbf{特点：}
\begin{itemize}
    \item 多谱线，需滤光片限制波长
    \item 强度均匀性需优于数个百分点
    \item 光源老化需定期校准（使用光谱曝光计）
\end{itemize}

\subsubsection{准分子激光（Excimer Laser）}

\textbf{KrF激光：}
\begin{itemize}
    \item \textbf{波长：}248 nm
    \item \textbf{用途：}深紫外（DUV）光刻
\end{itemize}

\textbf{ArF激光：}
\begin{itemize}
    \item \textbf{波长：}193 nm
    \item \textbf{用途：}当前主流光刻技术
    \item \textbf{浸没式（Immersion）：}通过水介质提高分辨率
\end{itemize}

\textbf{F$_2$激光：}
\begin{itemize}
    \item \textbf{波长：}157 nm
    \item \textbf{用途：}超准分子激光光刻
\end{itemize}

\textbf{脉冲特性：}
\begin{itemize}
    \item 脉冲宽度：20 nsec/脉冲
    \item 脉冲数：约50次脉冲
    \item 重复频率：几kHz
\end{itemize}

\subsection{曝光系统（Exposure Systems）}

\subsubsection{接触式曝光（Contact Printing）}

\textbf{原理：}掩模直接接触光刻胶表面

\textbf{优点：}
\begin{itemize}
    \item 简单
    \item 分辨率高
\end{itemize}

\textbf{缺点：}
\begin{itemize}
    \item 掩模易损坏
    \item 污染问题严重
    \item 1:1复制（无缩小）
\end{itemize}

\subsubsection{接近式曝光（Proximity Printing）}

\textbf{原理：}掩模与晶圆保持小间隙g

\textbf{衍射理论：}Fresnel衍射

\textbf{适用间隙范围：}
\begin{equation}
g \ll \frac{W^2}{\lambda}
\end{equation}

\textbf{最小分辨尺寸：}
\begin{equation}
W_{min} = \sqrt{g\lambda}
\end{equation}

\textbf{示例计算：}
\begin{itemize}
    \item $\lambda$ = 250 nm，g = 50 $\mu$m：$W_{min}$ = 4.5 $\mu$m
    \item $\lambda$ = 250 nm，g = 2 $\mu$m：$W_{min}$ = 0.9 $\mu$m
    \item $\lambda$ = 400 nm，g = 50 $\mu$m：$W_{min}$ = 3.5 $\mu$m
    \item $\lambda$ = 400 nm，g = 2 $\mu$m：$W_{min}$ = 0.7 $\mu$m
\end{itemize}

\textbf{优点：}
\begin{itemize}
    \item 简单
    \item 避免掩模损坏
\end{itemize}

\textbf{缺点：}
\begin{itemize}
    \item 分辨率受间隙限制
    \item 1:1复制
\end{itemize}

\subsubsection{投影式曝光（Projection Printing）}

\textbf{原理：}通过透镜系统将掩模图形缩小投影到晶圆

\textbf{缩小倍数（Demagnification）：}
\begin{itemize}
    \item 4X（最常用）
    \item 5X
    \item 10X
\end{itemize}

\textbf{优点：}
\begin{itemize}
    \item 超高分辨率（使用深紫外光）
    \item 掩模图形4倍放大，制作容易
    \item 掩模受保护，不接触晶圆
\end{itemize}

\textbf{缺点：}
\begin{itemize}
    \item 光学系统复杂
    \item 设备昂贵
\end{itemize}

\subsubsection{对比：接近式 vs 投影式}

\begin{longtable}{|L{5cm}|L{5cm}|L{5cm}|}
\hline
\textbf{特性} & \textbf{接近式} & \textbf{投影式} \\
\hline
\endfirsthead
\hline
\textbf{特性} & \textbf{接近式} & \textbf{投影式} \\
\hline
\endhead
\hline
\endfoot

复杂度 & 简单 & 复杂 \\
\hline

分辨率控制因素 & $\lambda$和间隙z & $\lambda$和数值孔径NA \\
\hline

掩模问题 & 1:1，可能损坏 & 4X，受保护 \\
\hline

成本 & 低 & 高 \\
\hline

\end{longtable}

\subsection{先进曝光设备}

\subsubsection{掩模对准机（Mask Aligner）}

\textbf{类型：}接近式印刷设备

\textbf{用途：}简单图形、教学实验

\subsubsection{步进机（Stepper）}

\textbf{工作方式：}
\begin{itemize}
    \item 逐个芯片曝光（Step and Repeat）
    \item 曝光一个die后，台移动到下一个位置
\end{itemize}

\textbf{特点：}
\begin{itemize}
    \item 高分辨率
    \item 精确对准
\end{itemize}

\subsubsection{步进扫描机（Step and Scan）}

\textbf{工作方式：}
\begin{itemize}
    \item 掩模和晶圆同时移动
    \item 狭缝扫描曝光
    \item 每个die扫描一次，然后步进到下一个die
\end{itemize}

\textbf{优势：}
\begin{itemize}
    \item 更大的曝光面积
    \item 更好的均匀性
    \item 主流生产设备（如ASML光刻机）
\end{itemize}

\subsection{光刻设备成本}

\textbf{成本趋势：}
\begin{itemize}
    \item 光学元件成本$\propto$(NA)$^3$
    \item 光学组件主导设备价格
    \item 历史成本增长：
    \begin{itemize}
        \item 1975-1985年：\$100,000
        \item 1985-2000年：\$1,000,000
        \item 2000-2010年：\$10,000,000
        \item 2010年后：接近\$100,000,000
    \end{itemize}
\end{itemize}

\subsection{分辨率增强技术（RET）}

近年来，临界尺寸显著小于光源波长，需要使用分辨率增强技术（tricks）：

\subsubsection{离轴照明（Off-Axis Illumination）}

改变照明角度，增强特定方向的分辨率

\subsubsection{相移掩模（Phase Shift Mask）}

利用相位差增强对比度

\subsubsection{光学邻近修正（OPC）}

Optical Proximity Correction：预先补偿邻近效应

\section{光刻材料}

\subsection{掩模（Mask）}

\subsubsection{掩模组成}

\textbf{1. 基板（Substrate）}

\textbf{要求：}
\begin{itemize}
    \item 高透光率（在曝光波长下）
    \item 低热膨胀系数
    \item 高平整度
    \item 最小的非线性
\end{itemize}

\textbf{常用材料：}
\begin{itemize}
    \item 石英（Quartz）
    \item 熔融硅（Fused silica）
    \item 硼硅酸盐玻璃（BSC）
\end{itemize}

\textbf{2. 吸收层（Absorber）}

\textbf{要求：}
\begin{itemize}
    \item 在曝光波长$\lambda$处无透射
    \item 与基板良好粘附
    \item 耐久性
\end{itemize}

\textbf{常用材料：}
\begin{itemize}
    \item 铬（Chrome）
    \item 乳剂（Emulsion）
    \item Fe$_2$O$_3$
\end{itemize}

\subsubsection{掩模极性（Mask Polarity）}

\textbf{暗场掩模（Dark Field）：}
\begin{itemize}
    \item 大部分区域不透光
    \item 需要曝光的区域透光
\end{itemize}

\textbf{明场掩模（Light Field）：}
\begin{itemize}
    \item 大部分区域透光
    \item 不需要曝光的区域不透光
\end{itemize}

\textbf{选择依据：}
\begin{itemize}
    \item 暗场：更少的邻近/背景曝光，缺陷影响小
    \item 明场：适合光栅等特定图形
\end{itemize}

\subsection{光刻胶（Resist）}

\subsubsection{光刻胶处理流程}

\textbf{完整流程：}
\begin{enumerate}
    \item \textbf{旋涂（Spincoat）：}光刻胶均匀涂布
    \item \textbf{涂胶后烘焙（Post-apply bake）：}溶剂挥发
    \item \textbf{曝光（Exposure）：}光化学反应
    \item \textbf{曝光后烘焙（Post-exposure bake, PEB）：}促进反应完成
    \item \textbf{显影（Development）：}选择性溶解
    \item \textbf{干燥（Drying）：}去除显影液
\end{enumerate}

\subsubsection{光刻胶厚度控制}

\textbf{影响因素：}
\begin{itemize}
    \item 旋涂速度（Spin speed）
    \item 光刻胶浓度（Resist concentration）
    \item 分子量（Molecular weight），影响粘度
\end{itemize}

\textbf{经验公式：}
\begin{equation}
t = K \cdot C^\beta \cdot \eta^\gamma \cdot \omega^{-\alpha}
\end{equation}

其中：
\begin{itemize}
    \item t = 厚度
    \item K = 校准常数
    \item C = 聚合物浓度（g/ml）
    \item $\eta$ = 本征粘度
    \item $\omega$ = 转速（rpm）
    \item $\beta$, $\gamma$, $\alpha$ = 经验因子
\end{itemize}

\subsubsection{光刻胶涂布要点}

\textbf{影响因素：}
\begin{itemize}
    \item 旋涂过程
    \item 粘度
    \item 旋涂速度
    \item 表面特性
    \item 环境（Direct ambient）
    \item 台阶覆盖（Step coverage）
    \item 粘附性（Adhesion）
\end{itemize}

\subsubsection{光刻胶要求}

\textbf{光学性能：}
\begin{itemize}
    \item \textbf{对比度（Contrast）：}区分曝光/未曝光区域
    \item \textbf{分辨率（Resolution）：}能形成的最小特征尺寸
    \item \textbf{灵敏度（Sensitivity）：}所需曝光剂量
\end{itemize}

\textbf{耐受性：}
\begin{itemize}
    \item 耐酸性
    \item 耐碱性
    \item 耐等离子体
    \item 耐离子轰击
    \item 耐反应离子刻蚀（RIE）
\end{itemize}

\textbf{工艺性能：}
\begin{itemize}
    \item 粘附性（Adhesion to substrates）
    \item 均匀性（Homogeneity）
    \item 易去除性（Ease of removal）
    \item 工艺可重复性（Process repeatability）
    \item 低缺陷密度（Low defect density）
    \item 台阶覆盖（Topography coverage）
    \item 温度稳定性（Temperature stability）
    \item 保质期（Shelf life）
    \item 纯净无污染（Purity from contaminants）
\end{itemize}

\subsubsection{光刻胶影响因素}

\textbf{材料因素：}
\begin{itemize}
    \item 玻璃化转变温度T$_g$
    \item 分子量
    \item 衬底原子序数Z
    \item 化学组成
\end{itemize}

\textbf{性能因素：}
\begin{itemize}
    \item 图形稳定性
    \item 分辨率
    \item 散射效应
    \item 刻蚀抗性
    \item 粘附性
\end{itemize}

\textbf{工艺因素：}
\begin{itemize}
    \item 显影（浓度、时间、温度）
    \item 烘焙（时间、温度）
    \item 后处理：去胶渣（descum）、天然氧化层去除等
\end{itemize}

\subsubsection{光刻胶类型}

\textbf{1. 负性交联型光刻胶（Negative Cross-linking Polymer）}

\textbf{原理：}曝光后聚合物交联，不溶于显影液

\textbf{特点：}
\begin{itemize}
    \item 曝光区域保留
    \item 未曝光区域溶解
\end{itemize}

\textbf{2. 正性酚醛树脂/重氮化合物光刻胶（Positive Novolak/Diazide）}

\textbf{中文名：}线型酚醛清漆

\textbf{原理：}曝光后光活性化合物变为可溶，增强溶解性

\textbf{特点：}
\begin{itemize}
    \item 曝光区域溶解
    \item 未曝光区域保留
    \item 分辨率高
\end{itemize}

\textbf{3. 正性链断裂型光刻胶（Positive Chain-Scission Polymer）}

\textbf{原理：}曝光导致聚合物链断裂，分子量降低，溶解性增强

\textbf{4. 自显影光刻胶（Self-Developing）}

\textbf{5. 化学放大型光刻胶（Chemically Amplified）}

\textbf{原理：}
\begin{itemize}
    \item 曝光产生少量酸或碱
    \item 后烘焙时催化链式反应
    \item 大幅提高灵敏度
\end{itemize}

\textbf{优势：}灵敏度极高，适合短波长光刻

\textbf{6. 硅化光刻胶（Silylated）}

\textbf{中文名：}甲烷硅基化的光刻胶

\subsubsection{光刻胶组成}

\textbf{1. 基体材料（Matrix Material）}

\textbf{也称：}树脂（Resin）

\textbf{功能：}
\begin{itemize}
    \item 作为粘结剂
    \item 粘附于衬底
    \item 对辐射不敏感
    \item 在显影液中快速溶解
    \item 物理或化学耐受性
\end{itemize}

\textbf{2. 敏化剂/光活性化合物（Sensitizer/Photoactive Compound）}

\textbf{功能：}
\begin{itemize}
    \item 无辐射时抑制溶解
    \item 辐射将其从溶解抑制剂转变为溶解增强剂
    \item 决定光刻胶的光化学响应
\end{itemize}

\textbf{作用机理：}
\begin{itemize}
    \item 溶解度变化可达约100:1
    \item 通过化学修饰改变溶解性
\end{itemize}

\textbf{3. 溶剂（Solvent）}

\textbf{功能：}
\begin{itemize}
    \item 保持光刻胶为液态
    \item 允许旋涂
    \item 浓度决定薄膜厚度
\end{itemize}

\subsubsection{光刻胶响应特性}

\textbf{非线性溶解：}

光刻胶的溶解不是线性的，存在阈值效应

\textbf{关键剂量：}
\begin{itemize}
    \item D$_0$ = 初始溶解剂量（Initial dissolution dose）
    \item D$_f$ = 100\%溶解剂量（100\% dissolution dose）
\end{itemize}

\textbf{理想情况：}
\begin{itemize}
    \item D$_f$ $\sim$ D$_0$ $\rightarrow$ 高对比度
    \item D$_f$小 $\rightarrow$ 高灵敏度
\end{itemize}

\subsubsection{对比度（Contrast）}

\textbf{定义：}
\begin{equation}
\gamma = [\log(D_f/D_0)]^{-1}
\end{equation}

\textbf{意义：}光刻胶区分未曝光和曝光区域的能力

\textbf{影响因素：}
\begin{itemize}
    \item 光刻胶化学性质
    \item 辐射类型
    \item 工艺条件
\end{itemize}

\textbf{要求：}$\gamma$值越大，对比度越高

\subsubsection{光刻胶灵敏度（Sensitivity）}

\textbf{本征灵敏度定义：}
\begin{equation}
\Phi = \frac{\text{化学反应数}}{\text{吸收光子数/电子数}}
\end{equation}

\textbf{特性：}
\begin{itemize}
    \item 光刻胶对辐射的响应
    \item 光谱响应：强吸收$\rightarrow$高灵敏度
    \item \textbf{权衡：}高灵敏度通常对应较差的分辨率
\end{itemize}

\subsubsection{分辨率限制因素}

\textbf{光刻胶与工艺：}

\textbf{本征因素：}
\begin{itemize}
    \item 聚合物尺寸（回转半径）
    \item 扩散
\end{itemize}

\textbf{工艺因素：}
\begin{itemize}
    \item 厚度
    \item 有效对比度
\end{itemize}

\section{光学光刻的关键因素}

\subsection{分辨率（Resolution）}

\subsubsection{光学分辨率理论}

\textbf{Fraunhofer衍射（远场投影）：}

当光通过小孔时，会发生衍射，形成Airy盘

\textbf{瑞利判据（Rayleigh Criterion）：}

\textbf{定义：}一个点源的中央最大值位于另一个点源Airy盘的第一极小值处

\textbf{瑞利公式：}
\begin{equation}
R = 0.61\frac{\lambda}{NA}
\end{equation}

其中：
\begin{itemize}
    \item R = 最小分辨尺寸
    \item $\lambda$ = 曝光波长
    \item NA = 数值孔径（Numerical Aperture）
\end{itemize}

\textbf{数值孔径定义：}
\begin{equation}
NA = n\sin(\alpha)
\end{equation}

其中：
\begin{itemize}
    \item n = 介质折射率（空气中为1，水中约为1.33）
    \item $\alpha$ = 半孔径角
\end{itemize}

\textbf{改进分辨率的方法：}
\begin{enumerate}
    \item 减小波长$\lambda$：使用深紫外、EUV光源
    \item 增大NA：
    \begin{itemize}
        \item 增大孔径角$\alpha$（受透镜设计限制）
        \item 增大折射率n（浸没式光刻）
    \end{itemize}
\end{enumerate}

\subsection{光学投影中的对比度}

\subsubsection{对比度定义}

\textbf{公式：}
\begin{equation}
\text{Contrast} = \frac{I_{max} - I_{min}}{I_{max} + I_{min}}
\end{equation}

其中：
\begin{itemize}
    \item I$_{max}$ = 最大光强
    \item I$_{min}$ = 最小光强
\end{itemize}

\subsubsection{调制传输函数（MTF）}

对比度即图形的调制传输函数（MTF）

\textbf{特性：}
\begin{itemize}
    \item 无衍射效应时：MTF = 1
    \item 有衍射效应时：MTF < 1
    \item 图形尺寸越小$\rightarrow$MTF越小
    \item 波长越短$\rightarrow$MTF越大
\end{itemize}

\subsection{临界尺寸（CD）控制}

\textbf{Critical Dimension Control影响因素：}

\subsubsection{曝光均匀性}

\begin{itemize}
    \item 工具依赖（Tool-dependent）
\end{itemize}

\subsubsection{照明}

\begin{itemize}
    \item 图形依赖（Pattern dependent）
    \item 邻近效应（Proximity effects）
\end{itemize}

\subsubsection{光刻胶}

\begin{itemize}
    \item 厚度（Thickness）
    \item 组成（Composition）
\end{itemize}

\subsubsection{显影液}

显影液浓度、温度、时间

\subsubsection{烘焙}

\begin{itemize}
    \item 曝光温度
    \item 显影温度
\end{itemize}

\subsubsection{图形转移}

刻蚀过程的精度

\subsection{焦深（Depth of Focus, DOF）}

\textbf{定义：}能够获得清晰成像的焦点深度范围

\textbf{重要性：}
\begin{itemize}
    \item 晶圆表面不完全平坦
    \item 需要足够的焦深容忍度
    \item DOF与分辨率存在权衡关系
\end{itemize}

\subsection{套刻精度（Placement Accuracy）}

\subsubsection{套刻类型}

\textbf{1. 层间套刻（Layer to Layer）}

不同掩模层之间的对准

\textbf{2. 特征间套刻（Feature to Feature）}

同一层内不同特征的对准

\subsubsection{写入方式}

\begin{itemize}
    \item \textbf{并行写入（Parallel writing）：}全晶圆同时曝光
    \item \textbf{串行写入（Serial writing）：}逐个区域曝光
\end{itemize}

\subsubsection{对准与套刻}

\textbf{对准标记（Alignment marks）：}
\begin{itemize}
    \item 在晶圆和掩模上预先制作的标记
    \item 用于精确定位
\end{itemize}

\textbf{套刻误差来源：}
\begin{itemize}
    \item 热膨胀/收缩
    \item 机械应力
    \item 台移精度
    \item 光学畸变
\end{itemize}

\section{光刻胶剖面与问题}

\subsection{光刻胶剖面类型}

光刻胶剖面取决于曝光和显影条件，可促进或阻碍后续工艺

\subsubsection{理想垂直剖面（Vertical）}

\textbf{特点：}侧壁垂直

\textbf{用途：}干法刻蚀

\subsubsection{底切剖面（Undercut）}

\textbf{特点：}底部比顶部窄

\textbf{用途：}剥离工艺（Lift-off）

\subsubsection{过切剖面（Overcut）}

\textbf{特点：}顶部比底部窄

\textbf{用途：}湿法刻蚀

\subsection{光刻胶污染}

\subsubsection{DUV光刻胶表面污染}

\textbf{问题：}深紫外（DUV）光刻胶顶部表面被碱暴露污染

\textbf{原因：}
\begin{itemize}
    \item 化学放大型光刻胶对碱敏感
    \item 环境中的碱性物质（如胺类）
\end{itemize}

\textbf{影响：}
\begin{itemize}
    \item 表面失活
    \item T-topping（顶部轮廓异常）
\end{itemize}

\textbf{解决方法：}
\begin{itemize}
    \item 控制环境
    \item 使用顶部抗反射涂层（TARC）
\end{itemize}

\subsection{其他光刻胶问题}

\subsubsection{胶渣（Scumming）}

\textbf{定义：}显影后应溶解区域仍有残留

\textbf{原因：}
\begin{itemize}
    \item 显影不充分
    \item 光刻胶曝光不足
\end{itemize}

\textbf{影响：}短路、图形失真

\subsubsection{图形坍塌（Pattern Collapse）}

\textbf{定义：}高深宽比的光刻胶线条倒塌

\textbf{原因：}
\begin{itemize}
    \item 表面张力（显影液干燥时）
    \item 光刻胶强度不足
    \item 深宽比过大
\end{itemize}

\textbf{解决方法：}
\begin{itemize}
    \item 优化光刻胶配方
    \item 改进干燥方法（如临界点干燥）
    \item 减小深宽比
\end{itemize}

\subsection{多层光刻胶（Multilayer Resists）}

\textbf{目的：}
\begin{itemize}
    \item 改善台阶覆盖
    \item 提高分辨率
    \item 减少基底反射
\end{itemize}

\textbf{结构：}
\begin{itemize}
    \item 底层：厚平坦化层
    \item 中间层：抗反射层或硬掩模
    \item 顶层：成像光刻胶层
\end{itemize}

\section{章节总结与知识点梳理}

\subsection{核心概念对照}

\begin{longtable}{|L{4.5cm}|L{10.5cm}|}
\hline
\textbf{概念} & \textbf{关键要点} \\
\hline
\endfirsthead
\hline
\textbf{概念} & \textbf{关键要点} \\
\hline
\endhead
\hline
\endfoot

光刻起源 & Alois Senefelder 1796年；化学印刷；lithos（石）+ graphy（写） \\
\hline

STM光刻 & 分辨率$\sim$0.3 nm；1原子/分钟；仅研究用 \\
\hline

现代光刻 & 5 nm EUV；193 nm ArF浸没式；化学放大胶；100片/小时 \\
\hline

光刻成本 & >60\%工艺成本；设备成本$\propto$(NA)$^3$；可达\$100M \\
\hline

分辨率 & 最小特征尺寸；尺寸控制精度 \\
\hline

对比度 & 区分曝光/未曝光区域能力 \\
\hline

产能 & 晶圆数/小时；受像素速率和胶灵敏度影响 \\
\hline

套刻精度 & 层间对准精度；包括层间和特征间 \\
\hline

Hg灯 & G线436 nm、H线405 nm、I线365 nm \\
\hline

准分子激光 & KrF 248 nm、ArF 193 nm、F$_2$ 157 nm \\
\hline

激光脉冲 & 20 nsec/脉冲；50次脉冲；几kHz重复率 \\
\hline

接触式曝光 & 掩模接触胶；简单但易损坏；1:1 \\
\hline

接近式曝光 & 间隙g；$W_{min}=\sqrt{g\lambda}$；Fresnel衍射 \\
\hline

投影式曝光 & 4X缩小；高分辨率；复杂昂贵；掩模受保护 \\
\hline

步进机 & Step and Repeat；逐die曝光 \\
\hline

步进扫描机 & Step and Scan；狭缝扫描；主流设备 \\
\hline

RET技术 & 离轴照明、相移掩模、光学邻近修正 \\
\hline

掩模基板 & 石英/熔融硅；高透光、低膨胀、高平整 \\
\hline

掩模吸收层 & 铬/Fe$_2$O$_3$；无透射、良好粘附、耐久 \\
\hline

暗场掩模 & 大部分不透光；少邻近曝光；缺陷影响小 \\
\hline

明场掩模 & 大部分透光；适合光栅 \\
\hline

光刻胶流程 & 旋涂$\rightarrow$前烘$\rightarrow$曝光$\rightarrow$PEB$\rightarrow$显影$\rightarrow$干燥 \\
\hline

胶厚控制 & $t \propto C^\beta \eta^\gamma \omega^{-\alpha}$ \\
\hline

负性胶 & 曝光后交联；曝光区保留 \\
\hline

正性胶 & 曝光后可溶；曝光区溶解；分辨率高 \\
\hline

化学放大胶 & 酸/碱催化链式反应；高灵敏度；用于短波长 \\
\hline

胶组成 & 基体（树脂）+ 敏化剂 + 溶剂 \\
\hline

敏化剂作用 & 溶解度变化100:1；抑制$\rightarrow$增强 \\
\hline

胶响应剂量 & D$_0$初始溶解；D$_f$完全溶解；理想D$_f\sim$D$_0$ \\
\hline

对比度公式 & $\gamma=[\log(D_f/D_0)]^{-1}$；$\gamma$越大越好 \\
\hline

胶灵敏度 & $\Phi=$化学反应数/光子数；高灵敏度$\leftrightarrow$低分辨率 \\
\hline

瑞利公式 & $R=0.61\lambda/NA$；$NA=n\sin(\alpha)$ \\
\hline

分辨率提高 & 减小$\lambda$；增大NA（大孔径、浸没） \\
\hline

光学对比度 & $(I_{max}-I_{min})/(I_{max}+I_{min})$；即MTF \\
\hline

MTF特性 & 无衍射MTF=1；尺寸小MTF小；波长短MTF大 \\
\hline

CD控制 & 曝光均匀性、照明、胶厚度、显影、烘焙、图形转移 \\
\hline

焦深DOF & 清晰成像的深度范围；与分辨率权衡 \\
\hline

垂直剖面 & 干法刻蚀用 \\
\hline

底切剖面 & 剥离工艺用 \\
\hline

过切剖面 & 湿法刻蚀用 \\
\hline

DUV胶污染 & 表面被碱污染；T-topping；需环境控制 \\
\hline

胶渣scumming & 显影不充分；曝光不足；导致短路 \\
\hline

图形坍塌 & 高深宽比；表面张力；需优化配方 \\
\hline

多层胶 & 底层平坦化、中层抗反射、顶层成像 \\
\hline

\end{longtable}

\subsection{重要公式和数据}

\begin{enumerate}
    \item \textbf{接近式曝光分辨率：}$W_{min} = \sqrt{g\lambda}$
    \item \textbf{瑞利判据：}$R = 0.61\lambda/NA$
    \item \textbf{数值孔径：}$NA = n\sin(\alpha)$
    \item \textbf{对比度：}$\gamma = [\log(D_f/D_0)]^{-1}$
    \item \textbf{光学对比度：}Contrast = $(I_{max}-I_{min})/(I_{max}+I_{min})$
    \item \textbf{胶厚度：}$t = K \cdot C^\beta \cdot \eta^\gamma \cdot \omega^{-\alpha}$
    \item \textbf{灵敏度：}$\Phi = \frac{\text{化学反应数}}{\text{吸收光子数}}$
    \item \textbf{光源波长：}I线365 nm、KrF 248 nm、ArF 193 nm、F$_2$ 157 nm、EUV 13.5 nm
    \item \textbf{缩小倍数：}4X、5X、10X
    \item \textbf{激光脉冲：}20 nsec/脉冲、约50脉冲、几kHz
    \item \textbf{产能：}100片/小时（300 mm晶圆）
    \item \textbf{光学成本：}Cost $\propto$ (NA)$^3$
\end{enumerate}

\subsection{关键工艺流程图}

\textbf{完整光刻工艺流程：}
\begin{enumerate}
    \item 表面清洗与准备
    \item 粘附促进剂处理（HMDS等）
    \item 旋涂光刻胶
    \item 前烘（90-120°C，60-90秒）
    \item 对准
    \item 曝光
    \item 曝光后烘焙PEB（化学放大胶）
    \item 显影
    \item 后烘（硬烘）
    \item 检查
    \item 图形转移（刻蚀或离子注入）
    \item 光刻胶去除
\end{enumerate}

\textbf{技术演进路线：}
\begin{enumerate}
    \item 接触式/接近式（早期）
    \item I线步进机（365 nm）
    \item DUV步进机（KrF 248 nm）
    \item DUV步进扫描机（ArF 193 nm）
    \item 浸没式光刻（193 nm + 水）
    \item EUV光刻（13.5 nm）
\end{enumerate}

\subsection{常见考点和易错点}

\begin{enumerate}
    \item \textbf{正性胶 vs 负性胶}
    \begin{itemize}
        \item 正性：曝光区溶解，分辨率高
        \item 负性：曝光区交联保留，膨胀问题
    \end{itemize}
    
    \item \textbf{瑞利公式理解}
    \begin{itemize}
        \item $R \propto \lambda/NA$
        \item 减小$\lambda$或增大NA都能提高分辨率
        \item NA受透镜和介质限制
    \end{itemize}
    
    \item \textbf{光刻胶三组分}
    \begin{itemize}
        \item 基体：粘结、溶解快
        \item 敏化剂：控制溶解性，100:1变化
        \item 溶剂：液态、旋涂、控制厚度
    \end{itemize}
    
    \item \textbf{高灵敏度与低分辨率的权衡}
    \begin{itemize}
        \item 强吸收$\rightarrow$高灵敏度
        \item 但通常导致分辨率下降
    \end{itemize}
    
    \item \textbf{MTF与图形尺寸}
    \begin{itemize}
        \item 无衍射：MTF = 1
        \item 有衍射：MTF < 1
        \item 尺寸越小，MTF越小
    \end{itemize}
    
    \item \textbf{对比度定义}
    \begin{itemize}
        \item 胶对比度：$\gamma = [\log(D_f/D_0)]^{-1}$
        \item 光学对比度：$(I_{max}-I_{min})/(I_{max}+I_{min})$
        \item 两者不同但相关
    \end{itemize}
    
    \item \textbf{掩模极性选择}
    \begin{itemize}
        \item 暗场：少背景曝光，缺陷影响小
        \item 选择取决于图形密度和类型
    \end{itemize}
    
    \item \textbf{化学放大胶特点}
    \begin{itemize}
        \item 需要PEB步骤
        \item 高灵敏度
        \item 对环境敏感（碱污染）
    \end{itemize}
\end{enumerate}

\subsection{与其他章节的联系}

\begin{itemize}
    \item \textbf{清洗（Chapter 2）：}光刻前需要清洁表面，光刻后需去除光刻胶
    \item \textbf{氧化：}氧化层可作为光刻掩模
    \item \textbf{刻蚀：}光刻胶作为刻蚀掩模，剖面影响刻蚀效果
    \item \textbf{离子注入：}光刻胶作为注入掩模
    \item \textbf{薄膜沉积：}台阶覆盖影响光刻胶涂布，剥离工艺需底切剖面
    \item \textbf{CMP：}表面平坦化改善光刻质量
\end{itemize}

\section{复习建议}

\subsection{重点掌握内容}

\begin{enumerate}
    \item 光刻的基本概念和历史发展
    \item 三种曝光系统（接触、接近、投影）的原理和对比
    \item 光源类型和波长（Hg灯I线、KrF、ArF）
    \item 掩模的组成（基板和吸收层）和极性
    \item 光刻胶的类型（正性、负性、化学放大）
    \item 光刻胶的组成（基体、敏化剂、溶剂）
    \item 光刻工艺流程（6步）
    \item 瑞利公式和分辨率影响因素
    \item 对比度的两种定义
    \item 光刻胶剖面类型及应用
    \item 常见光刻问题（污染、胶渣、坍塌）
\end{enumerate}

\subsection{理解性内容}

\begin{enumerate}
    \item 为什么现代光刻成本如此高？（光学元件复杂性）
    \item 为什么投影式曝光成为主流？（高分辨率、掩模保护）
    \item 如何提高光学光刻分辨率？（短波长、大NA、RET）
    \item 化学放大胶为什么需要PEB？（催化反应完成）
    \item 为什么高灵敏度与高分辨率难以兼得？（吸收与散射）
    \item MTF如何影响图形质量？（衍射导致对比度下降）
    \item 为什么需要多层光刻胶？（台阶覆盖、抗反射）
    \item 步进扫描机相比步进机的优势？（更大曝光面积、更好均匀性）
\end{enumerate}

\subsection{计算题准备}

\begin{enumerate}
    \item 给定$\lambda$和间隙g，计算接近式曝光的$W_{min}$
    \item 给定$\lambda$和NA，计算投影光刻分辨率R
    \item 给定n和$\alpha$，计算NA
    \item 给定D$_0$和D$_f$，计算对比度$\gamma$
    \item 给定I$_{max}$和I$_{min}$，计算光学对比度
    \item 浸没式光刻分辨率改善计算（水的n$\approx$1.33）
\end{enumerate}

\subsection{记忆技巧}

\begin{itemize}
    \item \textbf{光刻流程：}涂$\rightarrow$烘$\rightarrow$光$\rightarrow$烘$\rightarrow$显$\rightarrow$干（6步，两次烘焙）
    \item \textbf{Hg灯GHI：}G(436) > H(405) > I(365)，波长递减
    \item \textbf{准分子激光：}KrF(248) > ArF(193) > F$_2$(157)
    \item \textbf{正负胶：}正性胶曝光区\textbf{正}好溶解
    \item \textbf{胶三组分：}基（体）+ 感（敏化剂）+ 溶（剂）
    \item \textbf{敏化剂作用：}无光抑制$\rightarrow$有光增强，100倍差异
    \item \textbf{分辨率公式：}R = 0.61$\lambda$/NA，记住系数0.61
    \item \textbf{NA提高：}大孔径or浸没（水n=1.33）
    \item \textbf{剖面应用：}垂直干刻、底切剥离、过切湿刻
\end{itemize}

\subsection{与实际应用的联系}

\begin{enumerate}
    \item \textbf{摩尔定律：}光刻技术的进步是推动因素
    \item \textbf{EUV光刻：}5 nm及以下节点的关键技术
    \item \textbf{多重图案化：}在单一曝光分辨率限制下获得更小特征
    \item \textbf{浸没式光刻：}通过水介质将193 nm光刻延伸至更小节点
    \item \textbf{计算光刻：}OPC等技术应对亚波长光刻挑战
    \item \textbf{成本考量：}光刻占IC制造成本60%以上，推动新技术平衡性能与成本
\end{enumerate}

\end{document}
